% TOP_PROBLEM: {"id": 2814, "title": "Kirby Problem 3.16", "category_id": 7, "category": {"id": 7, "name": "topology", "display_name": "Topology", "description": "Properties preserved under continuous deformations.", "slug": "topology", "order_index": 7, "created_at": "2026-07-31T15:26:25.671Z"}, "status": "open", "problem_number": "KP-3.16", "set_id": 11}
% FIRST_POSTED: 2026-10-01
% ENTRY_KIND: statement_only
% UnsolvedMath ID 2814; source catalogue code KP-3.16
% !TeX program = lualatex
% Definition and sources only; this file is not an LLM solution attempt.
\documentclass[11pt,a4paper]{article}
\usepackage[margin=25mm]{geometry}
\usepackage{fontspec}
\setmainfont{lmroman10-regular.otf}[BoldFont=lmroman10-bold.otf,
 ItalicFont=lmroman10-italic.otf,BoldItalicFont=lmroman10-bolditalic.otf]
\usepackage{amsmath,amssymb,mathtools}
\usepackage{unicode-math}
\setmathfont{latinmodern-math.otf}
\AtBeginDocument{\let\setminus\smallsetminus}
\usepackage{xurl,hyperref}
\hypersetup{colorlinks=true,urlcolor=blue}
\setcounter{secnumdepth}{0}
\setlength{\parindent}{0pt}
\setlength{\parskip}{5pt}
\setlength{\emergencystretch}{3em}
\begin{document}
\section*{Kirby Problem 3.16}\label{2814}
{\small UnsolvedMath ID 2814 · KP-3.16 · Topology · Kirby's Problems in Low-Dimensional Topology}
\subsection{Definitions and mathematical statement}
Let $M$ be a connected complete finite-volume hyperbolic three-manifold.
A closed geodesic is a periodic geodesic $\gamma:\mathbb R/L\mathbb Z\to M$
parametrized by arclength. It is simple when this map is an embedding.
We count geodesics by their geometric images, identifying reparametrizations
and orientation reversal. Traversing one geodesic multiple times therefore
does not create new simple geodesics.

Is the set of simple closed geodesics in every such $M$ infinite? Equivalently,
for every positive integer $N$, must $M$ contain $N$ distinct geometric
images of embedded closed geodesics? No requirement is made that these
different geodesics be disjoint from one another.

The source records an affirmative answer for cusped manifolds and asks for
the general statement, in particular for all closed hyperbolic three-manifolds.
A shortest closed geodesic supplies one simple example in a closed manifold,
but this does not imply infinitude. Similarly, infinitely many free homotopy
classes of loops do not suffice: their geodesic representatives must be
shown to avoid self-intersections. The metric is the given hyperbolic metric,
not a perturbation selected to improve the behavior of geodesics.
\subsection{Short English statement}
Does every finite-volume hyperbolic three-manifold have infinitely many distinct embedded closed geodesics, including in the closed case?
\subsection{Sources}
\begin{itemize}
\item[S1] ULAM.AI and the UnsolvedMath Contributors, supplied problems.json, record 2814 (KP-3.16), Kirby's Problems in Low-Dimensional Topology. Catalogue identity and source transcription; CC BY 4.0.
\url{https://huggingface.co/datasets/ulamai/UnsolvedMath}.
\item[S2] K3: A New Problem List in Low-Dimensional Topology, Problem 3.16. Expanded formulation of the supplied catalogue statement.
\url{https://math.berkeley.edu/sites/default/files/surv-295-ruberman-watermarked-author-pdf.pdf}.
\end{itemize}
\end{document}
